\setauthor{Sebastian Radić \& Luis Schörgendorfer}
\chapter{Grenzen und zukünftige Arbeiten}
\label{chap:grenzen_zukunft}

\section{Bekannte Einschränkungen}
Der \textit{SmartBillConverter} funktioniert -- das belegen die validierten Testergebnisse.
Trotzdem gibt es Grenzen, die im Rahmen dieser Diplomarbeit nicht überwunden wurden
und die beim produktiven Einsatz bekannt sein sollten.

\textbf{OCR-Qualität bei schlechten Vorlagen.}
Die OCR-Pipeline funktioniert gut bei sauber gescannten oder fotografierten Rechnungen.
Bei handgeschriebenen Rechnungen, stark verzerrten Bildern oder sehr niedrig aufgelösten
Scans nimmt die Erkennungsqualität deutlich ab. Tesseract liefert in diesen Fällen
fehlerhafte Texte, die auch durch die KI-Normalisierung nur bedingt korrigiert werden
können, da die Eingabegrundlage für das Sprachmodell schlicht zu schlecht ist.

\textbf{Abhängigkeit von der externen Gemini~API.}
Die KI-Normalisierung setzt voraus, dass die Google Gemini~API erreichbar ist.
Das bedeutet, dass das System ohne Internetverbindung nicht funktioniert.
Wichtiger noch: Rechnungsdaten werden an einen externen Dienst übertragen,
was je nach Unternehmensrichtlinien datenschutzrechtlich problematisch sein kann.
Im Rahmen dieser Arbeit wurde das als vertretbar eingestuft, da der Einsatz
rein intern ist. Für einen breiteren Betrieb wäre das neu zu bewerten.

\textbf{Kein produktionsreifes Deployment.}
Das System läuft über eine lokale Docker-Umgebung. Es gibt keine Absicherung
für den Betrieb auf einem öffentlich erreichbaren Server: kein HTTPS-Termination,
keine Secrets-Verwaltung, kein Monitoring. Das war bewusst außerhalb des
Projektscopes, schränkt aber die unmittelbare Einsetzbarkeit ein.

\textbf{Begrenzte Standardunterstützung.}
Das System unterstützt ausschließlich ebInterface~6.1 und ZUGFeRD~2.3,
da das die beiden Formate sind, die die PROGRAMMIERFABRIK GmbH benötigt.
Weitere in Europa gebräuchliche Standards -- etwa XRechnung (Deutschland)
oder Peppol~BIS -- sind nicht implementiert.

\textbf{KI-generiertes XML erfordert manuelle Überprüfung.}
Auch wenn alle Testrechnungen die offiziellen Validatoren bestanden haben,
bedeutet eine erfolgreiche XSD-Validierung nicht automatisch, dass die Inhalte
semantisch korrekt sind. Kleine Fehler -- ein falscher Betrag, ein nicht
erkannter Empfängername -- können auftreten und müssen vom Benutzer vor der
Weiterverwendung geprüft werden.

\section{Zukünftige Erweiterungen}
Der \textit{SmartBillConverter} ist als Prototyp konzipiert, hat aber eine klare
Basis, auf der weitergebaut werden kann.

\textbf{Weitere E-Rechnungsstandards.}
Die Verarbeitungspipeline ist generisch aufgebaut: Textextraktion,
KI-Normalisierung und XML-Generierung sind voneinander entkoppelt.
Das Hinzufügen eines neuen Zielformats bedeutet im Wesentlichen,
einen neuen \texttt{XmlService} zu implementieren und das entsprechende
Schema zu kennen. XRechnung (basierend auf UBL und CII) und Peppol~BIS
wären naheliegende nächste Schritte, besonders wenn das Tool auch für
den deutschen Markt relevant werden soll.

\textbf{Lokales KI-Modell (On-Premise).}
Die Abhängigkeit von der Gemini~API ließe sich durch ein lokal betriebenes
Sprachmodell lösen. Modelle wie LLaMA oder Mistral, die über Ollama lokal
ausgeführt werden können, wären ein möglicher Ersatz.\footnote{Vgl. Meta AI: \textit{Llama -- Open Foundation and Fine-Tuned Chat Models}, \url{https://llama.meta.com/}, letzter Zugriff am 20.01.2026}
Das würde Datenschutzbedenken entschärfen und den Offline-Betrieb ermöglichen,
geht aber mit höheren Hardware-Anforderungen einher.

\textbf{Batch-Verarbeitung und automatischer Import.}
Derzeit müssen Rechnungen manuell über das Frontend hochgeladen werden.
Eine automatische Verarbeitung -- zum Beispiel durch Überwachung eines lokalen
Ordners oder eines E-Mail-Postfachs -- würde den manuellen Aufwand weiter
reduzieren und den praktischen Nutzen für den Tagesbetrieb deutlich erhöhen.

\textbf{Feedback-Loop und aktives Lernen.}
Wenn ein Benutzer einen KI-Fehler im generierten XML korrigiert, könnte
diese Korrektur als Trainingsdatum verwendet werden, um das System über
die Zeit zu verbessern. Ein einfacher Ansatz wäre, korrigierte Rechnungspaare
(Rohtext +\ korrektes JSON) zu sammeln und für Few-Shot-Prompting oder
Fine-Tuning zu nutzen.
